\documentclass[12pt,a4paper]{article}
\usepackage[utf8]{inputenc}
\usepackage[T5]{fontenc}
\usepackage[vietnamese]{babel}
\usepackage[margin=1.5cm,top=1.4cm,bottom=1.4cm]{geometry}
\usepackage{amsmath,amssymb,array,tabularx,enumitem,fancyhdr,tikz,xcolor}
\usetikzlibrary{calc,angles,quotes,arrows.meta,babel,patterns}
\definecolor{ink}{RGB}{30,55,70}

\newcommand{\headerQuote}{%
  \ifcase\value{page}\relax
    \or Math gives us hope that every problem has a solution% Trang 1
    \or Math + You = Champion% Trang 2
    \or Where Frustration Becomes Inspiration% Trang 3
    \or Turning Confusion into Clarity% Trang 4
    \or Don't be an X, be a Y of mathematics% Trang 5
    \or Don't get mad. Learn to add% Trang 6
    \or Math gives us hope that every problem has a solution% Trang 7
    \or Math + You = Champion% Trang 8
    \or Where Frustration Becomes Inspiration% Trang 9
    \or Turning Confusion into Clarity% Trang 10
    \or Don't be an X, be a Y of mathematics% Trang 11
    \or Don't get mad. Learn to add% Trang 12
    \else Math gives us hope that every problem has a solution%
  \fi
}
\pagestyle{fancy}\fancyhf{}
\lhead{\small\textit{\headerQuote}}
\rhead{\small\textit{Người biên soạn: Dương Gia Bảo}}
\rfoot{\small Trang \thepage}
\renewcommand{\headrulewidth}{0.4pt}
\setlength{\headheight}{18pt}
\setlength{\parindent}{0pt}
\setlength{\parskip}{1.5pt}
\setlength{\tabcolsep}{4pt}
\renewcommand{\arraystretch}{1.12}
\newcommand{\parttitle}[1]{\par\vspace{2pt}{\large\bfseries\color{ink}#1}\par\vspace{1pt}}
\newcommand{\four}[4]{\par\noindent\begin{tabularx}{\linewidth}{@{}*{4}{>{\raggedright\arraybackslash}X}@{}}\textbf{A.} #1&\textbf{B.} #2&\textbf{C.} #3&\textbf{D.} #4\end{tabularx}}
\newcommand{\two}[4]{\par\noindent\begin{tabularx}{\linewidth}{@{}*{2}{>{\raggedright\arraybackslash}X}@{}}\textbf{A.} #1&\textbf{B.} #2\\[2pt]\textbf{C.} #3&\textbf{D.} #4\end{tabularx}}
\newcommand{\answerbox}{\hfill\textit{Đáp số:}\ \fbox{\rule{0pt}{12pt}\hspace{2.2cm}}}
\setlist[enumerate,1]{label=\textbf{Câu \arabic*.},leftmargin=1.4cm,labelsep=.15cm,itemsep=1.5pt,topsep=1.0pt,parsep=0pt}
\setlist[enumerate,2]{label=\alph*),leftmargin=.6cm,itemsep=1.0pt,topsep=1.0pt,parsep=0pt}

\begin{document}
\noindent
\begin{minipage}[c]{0.38\textwidth}
    \centering
    \textbf{LÀM TOÁN ĐỂ HẠNH PHÚC}\\[-2pt]
    \rule{3.2cm}{0.6pt}\\[4pt]
    \textbf{ĐỀ CHÍNH THỨC}
\end{minipage}\hfill
\begin{minipage}[c]{0.35\textwidth}
    \centering
    \textbf{ĐỀ ÔN THI GK1}\\
    \textbf{NĂM HỌC 2026 -- 2027}\\[2pt]
    \textbf{ĐỀ SỐ 06}\\
    \textbf{Môn: TOÁN -- Khối lớp 10}\\
    \textit{\small Thời gian làm bài: 90 phút\\(không kể thời gian phát đề)}
\end{minipage}

\vspace{0.2cm}
\noindent
\begin{minipage}[b]{0.78\textwidth}
    \textbf{Họ và tên:}\dotfill\textbf{SBD:}\makebox[3.5cm]{\dotfill}
\end{minipage}\hfill
\begin{minipage}[b]{0.20\textwidth}
    \hfill
    \setlength{\fboxrule}{0.8pt}%
    \fbox{\begin{minipage}{2.3cm}
        \centering
        \textbf{\small Mã đề thi}\\
        \textbf{0006}
    \end{minipage}}
\end{minipage}
\vspace{0.2cm}

\parttitle{PHẦN I. Trắc nghiệm nhiều lựa chọn}
\textit{Thí sinh trả lời từ câu 1 đến câu 12. Mỗi câu hỏi chỉ chọn một phương án đúng duy nhất.}
\begin{enumerate}
\item Cho tập hợp $A = \{x \in \mathbb{N} \mid (2x - 4)(x^2 - 9) = 0\}$. Viết tập hợp $A$ bằng cách liệt kê các phần tử.
\four{$A = \{2\}$.}{$A = \{2; 3\}$.}{$A = \{-3; 2; 3\}$.}{$A = \{3\}$.}

\item Cho hai tập hợp số $A = (-3; 4]$ và $B = [1; 6)$. Tập hợp $A \cup B$ bằng
\four{$[1; 4]$.}{$(-3; 1)$.}{$(-3; 6)$.}{$[-3; 6]$.}

\item Tập hợp gồm tất cả các số thực $x$ thỏa mãn $-2 \leq x < 5$ được viết dưới dạng khoảng, đoạn, nửa khoảng là
\four{$[-2; 5)$.}{$(-2; 5]$.}{$(-2; 5)$.}{$[-2; 5]$.}

\item Cho tập hợp $M$ có đúng $4$ phần tử. Số tập hợp con của tập hợp $M$ là
\four{$4$.}{$8$.}{$12$.}{$16$.}

\item Bất phương trình nào sau đây là bất phương trình bậc nhất hai ẩn $x, y$?
\four{$2x^2 + y \leq 3$.}{$x - 2y^2 > 0$.}{$xy + 3y \geq 1$.}{$3x - 4y + 2 \leq 0$.}

\item Điểm nào sau đây thuộc miền nghiệm của hệ bất phương trình $\begin{cases} x - y \geq 0 \\ 2x + y \leq 6 \\ y \geq 0 \end{cases}$?
\four{$(1; 3)$.}{$(-1; 2)$.}{$(2; 1)$.}{$(3; 2)$.}

\item Cặp số $(x; y)$ nào sau đây là một nghiệm của bất phương trình $3x - 2y + 6 > 0$?
\four{$(0; 0)$.}{$(-3; 0)$.}{$(-2; 1)$.}{$(0; 4)$.}

\item Hệ bất phương trình nào sau đây là hệ bất phương trình bậc nhất hai ẩn $x, y$?
\two{$\begin{cases} x^2 - y \leq 1 \\ 2x + y \geq 0 \end{cases}$}{$\begin{cases} 2x - y \geq 1 \\ x + 3y \leq 4 \end{cases}$}{$\begin{cases} xy \leq 2 \\ x - y > 3 \end{cases}$}{$\begin{cases} x - y^2 \leq 5 \\ y \geq 0 \end{cases}$}

\item Cho góc nhọn $\alpha$ thỏa mãn $\sin\alpha = \dfrac{3}{5}$. Giá trị của biểu thức $P = 2\cos\alpha + \tan\alpha$ bằng
\four{$\dfrac{47}{20}$.}{$\dfrac{11}{5}$.}{$\dfrac{31}{20}$.}{$\dfrac{23}{10}$.}

\item Cho tam giác $ABC$ có nửa chu vi $p = \dfrac{a+b+c}{2}$ và bán kính đường tròn nội tiếp là $r$. Diện tích $S$ của tam giác $ABC$ được tính theo công thức nào sau đây?
\four{$S = 2pr$.}{$S = pr$.}{$S = \dfrac{1}{2}pr$.}{$S = \dfrac{p}{r}$.}

\item Cho tam giác $ABC$ có các cạnh $BC = a, CA = b, AB = c$. Khẳng định nào sau đây là đúng?
\two{$\cos A = \dfrac{b^2 + c^2 - a^2}{bc}$.}{$\cos A = \dfrac{a^2 + c^2 - b^2}{2ac}$.}{$\cos A = \dfrac{b^2 + c^2 - a^2}{2bc}$.}{$\cos A = \dfrac{b^2 + c^2 + a^2}{2bc}$.}

\item Giá trị của $\cos 120^\circ$ bằng
\four{$\dfrac{1}{2}$.}{$\dfrac{\sqrt{3}}{2}$.}{$-\dfrac{\sqrt{3}}{2}$.}{$-\dfrac{1}{2}$.}
\end{enumerate}

\parttitle{PHẦN II. Trắc nghiệm đúng/sai}
\textit{Thí sinh trả lời câu 1 và câu 2. Trong mỗi ý a), b), c), d) ở mỗi câu, thí sinh chọn đúng hoặc sai.}
\begin{enumerate}
\item Một xưởng may gia công nhận sản xuất hai loại áo: áo sơ mi (loại A) và áo khoác gió (loại B).
\begin{itemize}[leftmargin=0.6cm, itemsep=1.5pt]
    \item Để may một chiếc áo loại A cần $2\text{ m}$ vải và $3\text{ giờ}$ công may.
    \item Để may một chiếc áo loại B cần $3\text{ m}$ vải và $2\text{ giờ}$ công may.
\end{itemize}
Xưởng may có sẵn $120\text{ m}$ vải và quỹ thời gian lao động tối đa là $100\text{ giờ}$. Gọi $x$ và $y$ ($x, y \in \mathbb{N}$) lần lượt là số lượng áo loại A và loại B mà xưởng may sản xuất.
\begin{enumerate}
\item Lượng vải cần dùng để may $x$ chiếc áo loại A và $y$ chiếc áo loại B là $2x + 3y\text{ (m)}$.
\item Nếu xưởng may sản xuất $20$ chiếc áo loại A thì số áo loại B tối đa có thể may được để không vượt quá lượng vải có sẵn là $26$ chiếc.
\item Hệ bất phương trình mô tả các điều kiện ràng buộc đối với $x$ và $y$ là:
\[
\begin{cases}
x \geq 0,\ y \geq 0 \\
2x + 3y \leq 120 \\
3x + 2y \leq 100
\end{cases}
\]
\item Xưởng may có thể sản xuất đồng thời $25$ chiếc áo loại A và $20$ chiếc áo loại B mà không vi phạm cả hai điều kiện về lượng vải và thời gian làm việc.
\end{enumerate}

\item \begin{minipage}[t]{0.60\linewidth}
Tại một trạm quan sát hải đăng đặt tại điểm $A$ bên bờ biển, nhân viên quan sát theo dõi hai tàu đánh cá neo đậu tại các vị trí $B$ và $C$. Qua thiết bị đo đạc, nhân viên xác định được:
\begin{itemize}[leftmargin=0.5cm, itemsep=1.5pt]
    \item Khoảng cách từ trạm $A$ đến tàu $B$ là $AB = 6\text{ km}$.
    \item Khoảng cách từ trạm $A$ đến tàu $C$ là $AC = 10\text{ km}$.
    \item Góc quan sát giữa hai tàu là $\widehat{BAC} = 60^\circ$.
\end{itemize}
\end{minipage}\hfill
\begin{minipage}[t]{0.38\linewidth}
\vspace{-0.3\baselineskip}
\centering
\begin{tikzpicture}[scale=0.38, font=\normalsize, >=Stealth, line join=round, line cap=round]
    \coordinate (A) at (0,0);
    \coordinate (B) at (6,0);
    \coordinate (C) at ({10*cos(60)}, {10*sin(60)});
    
    % Tam giác ABC
    \draw[line width=1.3pt, ink] (A) -- (B) node[midway, below=2pt] {$6\text{ km}$} -- (C) -- cycle;
    \draw[line width=1.1pt, teal!80!black] (A) -- (C) node[midway, above left=2pt] {$10\text{ km}$};
    
    % Góc BAC
    \draw pic[draw, angle radius=0.8cm, "$60^\circ$", angle eccentricity=1.5, font=\small\bfseries] {angle = B--A--C};
    
    % Các điểm
    \fill (A) circle (3.5pt) node[below left=2pt, font=\normalsize\bfseries] {$A$};
    \fill (B) circle (3.5pt) node[below right=2pt, font=\normalsize\bfseries] {$B$};
    \fill (C) circle (3.5pt) node[above=2pt, font=\normalsize\bfseries] {$C$};
    
    % Đường cao AH
    \coordinate (H) at ($(B)!(A)!(C)$);
    \draw[dashed, line width=1.0pt, gray!80!black] (A) -- (H);
    \fill (H) circle (2.5pt) node[above right=1pt, font=\small] {$H$};
    \draw pic[draw, angle radius=0.3cm] {right angle = A--H--C};
\end{tikzpicture}
\end{minipage}
\begin{enumerate}
\item Giá trị của $\cos\widehat{BAC} = \cos 60^\circ = \dfrac{1}{2}$.
\item Diện tích tam giác $ABC$ tạo bởi trạm quan sát và hai tàu đánh cá bằng $15\sqrt{3}\text{ km}^2$.
\item Khoảng cách giữa hai tàu đánh cá $B$ và $C$ bằng $2\sqrt{19}\text{ km}$.
\item Một ca nô cứu hộ xuất phát từ trạm $A$ di chuyển thẳng ra tuyến đường nối hai tàu $B, C$ theo phương vuông góc với $BC$. Khoảng cách di chuyển ngắn nhất từ trạm $A$ đến tuyến đường $BC$ lớn hơn $6\text{ km}$.
\end{enumerate}
\end{enumerate}

\parttitle{PHẦN III. Trả lời ngắn}
\textit{Thí sinh trả lời từ câu 1 đến câu 4. Điền kết quả vào ô tương ứng.}
\begin{enumerate}
\item Cho hai tập hợp số $A = [-3; 7)$ và $B = (1; 10]$. Gọi $C = A \cap B$. Hỏi trong tập hợp $C$ có bao nhiêu số nguyên?
\answerbox

\item Cho miền nghiệm của hệ bất phương trình $\begin{cases} x \geq 0 \\ y \geq 0 \\ x + y \leq 6 \\ 2x + y \leq 8 \end{cases}$ là một miền tứ giác $D$ trên mặt phẳng tọa độ $Oxy$. Tìm giá trị lớn nhất của biểu thức $F(x, y) = 3x + 2y$ khi cặp số $(x; y)$ thuộc miền $D$.
\answerbox

\item Cho tam giác $ABC$ có độ dài ba cạnh lần lượt là $a = 13$, $b = 14$ và $c = 15$. Tính bán kính $r$ của đường tròn nội tiếp tam giác $ABC$.
\answerbox

\item Cho tam giác $ABC$ có độ dài ba cạnh $a, b, c$ thỏa mãn hệ thức:
\[
b(b^2 - a^2) = c(a^2 - c^2) \quad \text{và} \quad b \neq c.
\]
Tính số đo của góc $A$ (đơn vị: độ).
\answerbox
\end{enumerate}

\parttitle{PHẦN IV. Tự luận}
\textit{Thí sinh trình bày lời giải chi tiết cho các câu hỏi sau.}
\begin{enumerate}
\item Cho hai tập hợp số $A = [-2; 4)$ và $B = (1; 6]$.
\begin{enumerate}
    \item Xác định các tập hợp $A \cup B$ và $A \setminus B$.
    \item Biểu diễn tập hợp $A \setminus B$ trên trục số.
\end{enumerate}

\item Cho tam giác $ABC$ có độ dài cạnh $AC = 8$, $AB = 5$ và góc $\widehat{BAC} = 60^\circ$.
\begin{enumerate}
    \item Áp dụng định lý côsin, hãy tính độ dài cạnh $BC$.
    \item Tính diện tích tam giác $ABC$ và độ dài đường cao $h_a$ kẻ từ đỉnh $A$ xuống cạnh $BC$.
\end{enumerate}

\item \begin{minipage}[t]{0.58\linewidth}
Để đo chiều cao của một tháp truyền hình $TH$ dựng vuông góc với mặt đất nằm ngang (với $H$ là chân tháp trên mặt đất), do vướng hồ nước sâu nên người ta không thể đến gần chân tháp để đo trực tiếp.
\par
Các kỹ sư chọn hai vị trí cọc ngắm $A$ và $B$ trên mặt đất sao cho ba điểm $A, B, H$ thẳng hàng và khoảng cách $AB = 30\text{ m}$. Bằng giác kế, người ta đo được các góc nâng nhìn lên đỉnh tháp $T$ như sau:
\[
\widehat{TAH} = 40^\circ \quad \text{và} \quad \widehat{TBH} = 25^\circ.
\]
Tính chiều cao $TH$ của tháp truyền hình (làm tròn kết quả đến chữ số thập phân thứ hai).
\end{minipage}\hfill
\begin{minipage}[t]{0.40\linewidth}
\vspace{-0.3\baselineskip}
\centering
\begin{tikzpicture}[scale=0.75, font=\normalsize, >=Stealth, line join=round, line cap=round]
    % Điểm
    \coordinate (H) at (4.5,0);
    \coordinate (A) at (1.2,0);
    \coordinate (B) at (-1.8,0);
    \coordinate (T) at (4.5,2.77);
    
    % Mặt đất
    \draw[line width=1.2pt] (-2.5,0) -- (5.5,0);
    \fill[pattern=north east lines, pattern color=gray!60] (-2.5,-0.3) rectangle (5.5,0);
    
    % Tháp TH
    \draw[line width=2.0pt, ink] (H) -- (T) node[midway, right=2pt, font=\normalsize\bfseries] {$h$};
    \fill (T) circle (3pt) node[above=2pt, font=\normalsize\bfseries] {$T$};
    \fill (H) circle (3pt) node[below=3pt, font=\normalsize\bfseries] {$H$};
    \draw pic[draw, angle radius=0.3cm] {right angle = T--H--A};
    
    % Đường ngắm
    \draw[line width=1.1pt, teal!80!black] (T) -- (A);
    \draw[line width=1.1pt, blue!80!black] (T) -- (B);
    
    % Điểm A, B
    \fill (A) circle (3pt) node[below=3pt, font=\normalsize\bfseries] {$A$};
    \fill (B) circle (3pt) node[below=3pt, font=\normalsize\bfseries] {$B$};
    
    % Góc nâng
    \draw pic[draw, line width=0.8pt, angle radius=0.65cm, "$40^\circ$", angle eccentricity=1.6, font=\small\bfseries] {angle = H--A--T};
    \draw pic[draw, line width=0.8pt, angle radius=1.0cm, "$25^\circ$", angle eccentricity=1.4, font=\small\bfseries] {angle = H--B--T};
    
    % Đoạn AB
    \draw[<->, line width=0.9pt] (-1.8,-0.6) -- (1.2,-0.6) node[midway, fill=white, inner sep=1.5pt, font=\small\bfseries] {$30\text{ m}$};
    \draw[dotted, line width=0.8pt] (B) -- (-1.8,-0.6);
    \draw[dotted, line width=0.8pt] (A) -- (1.2,-0.6);
\end{tikzpicture}
\end{minipage}

\item Bác Nam có một trang trại và dự định đầu tư nuôi hai loại gia súc: bò sữa (loại I) và dê sinh sản (loại II).
\begin{itemize}[leftmargin=0.6cm, itemsep=1.5pt]
    \item Mỗi con bò sữa cần $3\text{ m}^2$ diện tích chuồng và $1\text{ triệu đồng}$ tiền vốn mua thức ăn trong một đợt nuôi.
    \item Mỗi con dê cần $1\text{ m}^2$ diện tích chuồng và $2\text{ triệu đồng}$ tiền vốn mua thức ăn trong một đợt nuôi.
\end{itemize}
Hiện tại trang trại có diện tích chuồng trại tối đa là $30\text{ m}^2$ và nguồn vốn dự trù cho thức ăn tối đa là $20\text{ triệu đồng}$. Qua kinh nghiệm thực tế, mỗi con bò sữa đem lại lợi nhuận $4\text{ triệu đồng}$, mỗi con dê đem lại lợi nhuận $3\text{ triệu đồng}$ sau mỗi đợt nuôi.
\begin{enumerate}
    \item Gọi $x$ và $y$ ($x, y \in \mathbb{N}$) lần lượt là số con bò sữa và số con dê bác Nam nuôi. Hãy thiết lập hệ bất phương trình mô tả các điều kiện ràng buộc đối với $x$ và $y$.
    \item Xác định tọa độ các đỉnh của miền nghiệm đa giác trên mặt phẳng tọa độ $Oxy$.
    \item Bác Nam nên nuôi bao nhiêu con bò sữa và bao nhiêu con dê để thu được lợi nhuận lớn nhất? Tính số tiền lợi nhuận lớn nhất đó.
\end{enumerate}
\end{enumerate}

\newpage
\section*{BẢNG ĐÁP ÁN TỔNG HỢP}

\subsection*{1. Phần I -- Trắc nghiệm nhiều lựa chọn}
\begin{center}
\setlength{\tabcolsep}{2.0pt}
\begin{tabular}{|c|c|c|c|c|c|c|c|c|c|c|c|}
\hline
\textbf{Câu 1} & \textbf{Câu 2} & \textbf{Câu 3} & \textbf{Câu 4} & \textbf{Câu 5} & \textbf{Câu 6} & \textbf{Câu 7} & \textbf{Câu 8} & \textbf{Câu 9} & \textbf{Câu 10} & \textbf{Câu 11} & \textbf{Câu 12} \\
\hline
\textbf{B} & \textbf{C} & \textbf{A} & \textbf{D} & \textbf{D} & \textbf{C} & \textbf{A} & \textbf{B} & \textbf{A} & \textbf{B} & \textbf{C} & \textbf{D} \\
\hline
\end{tabular}
\end{center}

\subsection*{2. Phần II -- Trắc nghiệm đúng/sai}
\begin{center}
\begin{tabular}{|c|c|c|c|c|}
\hline
\textbf{Câu} & \textbf{Ý a)} & \textbf{Ý b)} & \textbf{Ý c)} & \textbf{Ý d)} \\
\hline
\textbf{Câu 1} & Đ & Đ & Đ & S \\
\hline
\textbf{Câu 2} & Đ & Đ & Đ & S \\
\hline
\end{tabular}
\end{center}

\subsection*{3. Phần III -- Trả lời ngắn}
\begin{center}
\begin{tabular}{|c|c|c|c|}
\hline
\textbf{Câu 1} & \textbf{Câu 2} & \textbf{Câu 3} & \textbf{Câu 4} \\
\hline
\textbf{5} & \textbf{14} & \textbf{4} & \textbf{60} \\
\hline
\end{tabular}
\end{center}

\subsection*{4. Phần IV -- Tóm tắt kết quả tự luận}
\begin{itemize}[leftmargin=0.6cm, itemsep=2pt]
    \item \textbf{Câu 1:} a) $A \cup B = [-2; 6]$, $A \setminus B = [-2; 1]$. b) Biểu diễn đoạn $[-2; 1]$ trên trục số.
    \item \textbf{Câu 2:} a) $BC = 7$. b) Diện tích $S = 10\sqrt{3}$, đường cao $h_a = \dfrac{20\sqrt{3}}{7}$.
    \item \textbf{Câu 3:} Chiều cao tháp truyền hình $h \approx 31{,}49\text{ m}$.
    \item \textbf{Câu 4:} Nuôi $8$ con bò sữa và $6$ con dê; lợi nhuận lớn nhất là $50\text{ triệu đồng}$.
\end{itemize}

\newpage
\section*{HƯỚNG DẪN GIẢI CHI TIẾT}

\subsection*{Phần I. Trắc nghiệm nhiều lựa chọn}
\begin{itemize}[leftmargin=0.6cm, itemsep=3pt]
\item \textbf{Câu 1 (Chọn B):} Phương trình $(2x - 4)(x^2 - 9) = 0 \iff \left[\begin{array}{l} 2x - 4 = 0 \\ x^2 - 9 = 0 \end{array}\right. \iff \left[\begin{array}{l} x = 2 \\ x = 3 \\ x = -3 \end{array}\right.$. Vì $x \in \mathbb{N}$ nên ta nhận $x = 2$ và $x = 3$. Vậy $A = \{2; 3\}$.
\item \textbf{Câu 2 (Chọn C):} Hợp của hai tập hợp $A = (-3; 4]$ và $B = [1; 6)$ là tập hợp các số thực thuộc $A$ hoặc thuộc $B$, do đó $A \cup B = (-3; 6)$.
\item \textbf{Câu 3 (Chọn A):} Tập hợp các số thực $x$ thỏa mãn $-2 \leq x < 5$ có đầu mút bên trái lấy dấu bằng (ngoặc vuông $[$) và đầu mút bên phải không lấy dấu bằng (ngoặc tròn $))$, ký hiệu là $[-2; 5)$.
\item \textbf{Câu 4 (Chọn D):} Tập hợp có $n$ phần tử thì có đúng $2^n$ tập hợp con. Với $n = 4$, số tập con là $2^4 = 16$.
\item \textbf{Câu 5 (Chọn D):} Bất phương trình bậc nhất hai ẩn có dạng tổng quát $ax + by + c \leq 0$ (hoặc $\geq, <, >$) với $a, b$ không đồng thời bằng $0$. Do đó $3x - 4y + 2 \leq 0$ là bất phương trình bậc nhất hai ẩn. Các phương án A, B, C chứa $x^2, y^2, xy$ nên không phải.
\item \textbf{Câu 6 (Chọn C):} Thay tọa độ $(2; 1)$ vào hệ: $\begin{cases} 2 - 1 = 1 \geq 0 \\ 2(2) + 1 = 5 \leq 6 \\ 1 \geq 0 \end{cases}$ (thỏa mãn cả 3 bất phương trình). Vậy điểm $(2; 1)$ thuộc miền nghiệm.
\item \textbf{Câu 7 (Chọn A):} Thay cặp $(0; 0)$ vào bất phương trình: $3(0) - 2(0) + 6 = 6 > 0$ (mệnh đề đúng). Vậy $(0; 0)$ là một nghiệm.
\item \textbf{Câu 8 (Chọn B):} Hệ $\begin{cases} 2x - y \geq 1 \\ x + 3y \leq 4 \end{cases}$ gồm hai bất phương trình bậc nhất hai ẩn $x, y$.
\item \textbf{Câu 9 (Chọn A):} Vì $0^\circ < \alpha < 90^\circ$ nên $\cos\alpha > 0$. Ta có $\cos\alpha = \sqrt{1 - \sin^2\alpha} = \sqrt{1 - \left(\dfrac{3}{5}\right)^2} = \dfrac{4}{5}$. Suy ra $\tan\alpha = \dfrac{\sin\alpha}{\cos\alpha} = \dfrac{3/5}{4/5} = \dfrac{3}{4}$. Khi đó $P = 2\left(\dfrac{4}{5}\right) + \dfrac{3}{4} = \dfrac{8}{5} + \dfrac{3}{4} = \dfrac{47}{20}$.
\item \textbf{Câu 10 (Chọn B):} Công thức tính diện tích tam giác theo bán kính đường tròn nội tiếp là $S = pr$.
\item \textbf{Câu 11 (Chọn C):} Theo hệ quả định lý côsin: $\cos A = \dfrac{b^2 + c^2 - a^2}{2bc}$.
\item \textbf{Câu 12 (Chọn D):} Góc $120^\circ$ bù với góc $60^\circ$, do đó $\cos 120^\circ = -\cos 60^\circ = -\dfrac{1}{2}$.
\end{itemize}

\subsection*{Phần II. Trắc nghiệm đúng/sai}
\begin{itemize}[leftmargin=0.6cm, itemsep=4pt]
\item \textbf{Câu 1:}
\begin{enumerate}
    \item \textbf{Đúng:} Một áo loại A cần $2\text{ m}$ vải, một áo loại B cần $3\text{ m}$ vải nên tổng lượng vải cần may là $2x + 3y\text{ (m)}$.
    \item \textbf{Đúng:} Khi $x = 20$, lượng vải tiêu hao là $2(20) + 3y \leq 120 \iff 3y \leq 80 \iff y \leq 26{,}67$. Vì $y \in \mathbb{N}$ nên số áo loại B tối đa là $26$ chiếc.
    \item \textbf{Đúng:} Ràng buộc lượng vải $2x + 3y \leq 120$, thời gian làm việc $3x + 2y \leq 100$ và số lượng áo $x \geq 0, y \geq 0$.
    \item \textbf{Sai:} Khi $x = 25, y = 20$, kiểm tra thời gian: $3(25) + 2(20) = 75 + 40 = 115 > 100$, vượt quá quỹ thời gian tối đa $100\text{ giờ}$ của xưởng may.
\end{enumerate}

\item \textbf{Câu 2:}
\begin{enumerate}
    \item \textbf{Đúng:} $\cos 60^\circ = \dfrac{1}{2}$.
    \item \textbf{Đúng:} Diện tích tam giác $S = \dfrac{1}{2}AB\cdot AC\cdot\sin 60^\circ = \dfrac{1}{2}\cdot 6\cdot 10\cdot\dfrac{\sqrt{3}}{2} = 15\sqrt{3}\text{ km}^2$.
    \item \textbf{Đúng:} Áp dụng định lý côsin trong tam giác $ABC$:
    \[
    BC^2 = AB^2 + AC^2 - 2AB\cdot AC\cdot\cos A = 6^2 + 10^2 - 2\cdot 6\cdot 10\cdot\dfrac{1}{2} = 36 + 100 - 60 = 76.
    \]
    Suy ra $BC = \sqrt{76} = 2\sqrt{19}\text{ km}$.
    \item \textbf{Sai:} Khoảng cách ngắn nhất từ $A$ đến đường thẳng $BC$ là độ dài đường cao $AH$:
    \[
    AH = \dfrac{2S}{BC} = \dfrac{2\cdot 15\sqrt{3}}{2\sqrt{19}} = \dfrac{15\sqrt{3}}{\sqrt{19}} = \dfrac{15\sqrt{57}}{19} \approx 5{,}96\text{ km} < 6\text{ km}.
    \]
    Khẳng định khoảng cách lớn hơn $6\text{ km}$ là sai.
\end{enumerate}
\end{itemize}

\subsection*{Phần III. Trả lời ngắn}
\begin{itemize}[leftmargin=0.6cm, itemsep=4pt]
\item \textbf{Câu 1 (Đáp số: 5):} Giao của hai tập hợp $A = [-3; 7)$ và $B = (1; 10]$ là $C = A \cap B = (1; 7)$. Các số nguyên thuộc khoảng $(1; 7)$ là $2, 3, 4, 5, 6$. Tổng cộng có đúng $5$ số nguyên.
\item \textbf{Câu 2 (Đáp số: 14):} Miền nghiệm tứ giác $D$ giới hạn bởi các đường thẳng $x = 0, y = 0, x + y = 6, 2x + y = 8$. Tọa độ các đỉnh của tứ giác là:
\begin{itemize}
    \item Gốc tọa độ $O(0; 0) \implies F(0; 0) = 0$.
    \item Điểm $A(4; 0)$ (giao của $2x + y = 8$ với $Ox$) $\implies F(4; 0) = 3(4) + 2(0) = 12$.
    \item Điểm $B(2; 4)$ (giao điểm của $x + y = 6$ và $2x + y = 8$) $\implies F(2; 4) = 3(2) + 2(4) = 14$.
    \item Điểm $C(0; 6)$ (giao của $x + y = 6$ với $Oy$) $\implies F(0; 6) = 3(0) + 2(6) = 12$.
\end{itemize}
So sánh các giá trị, ta được $\max F(x, y) = 14$ đạt tại điểm $B(2; 4)$.
\item \textbf{Câu 3 (Đáp số: 4):} Nửa chu vi tam giác $p = \dfrac{13 + 14 + 15}{2} = 21$.
Áp dụng công thức Heron:
\[
S = \sqrt{p(p-a)(p-b)(p-c)} = \sqrt{21\cdot 8\cdot 7\cdot 6} = \sqrt{7056} = 84.
\]
Bán kính đường tròn nội tiếp là $r = \dfrac{S}{p} = \dfrac{84}{21} = 4$.
\item \textbf{Câu 4 (Đáp số: 60):} Biến đổi hệ thức đã cho:
\[
b(b^2 - a^2) = c(a^2 - c^2) \iff b^3 - ba^2 = ca^2 - c^3 \iff b^3 + c^3 = a^2(b + c).
\]
Áp dụng hằng đẳng thức tổng hai lập phương:
\[
(b + c)(b^2 - bc + c^2) = a^2(b + c).
\]
Vì $b, c$ là độ dài các cạnh của tam giác nên $b + c > 0$. Chia cả hai vế cho $b + c$, ta được:
\[
b^2 - bc + c^2 = a^2 \iff b^2 + c^2 - a^2 = bc.
\]
Theo định lý côsin: $\cos A = \dfrac{b^2 + c^2 - a^2}{2bc} = \dfrac{bc}{2bc} = \dfrac{1}{2}$.
Vì $0^\circ < \widehat{A} < 180^\circ$ nên $\widehat{A} = 60^\circ$.
\end{itemize}

\subsection*{Phần IV. Tự luận}
\begin{itemize}[leftmargin=0.6cm, itemsep=4pt]
\item \textbf{Câu 1 (0,5 điểm):}
\begin{enumerate}
    \item Ta có:
    \begin{itemize}
        \item $A \cup B = [-2; 4) \cup (1; 6] = [-2; 6]$.
        \item $A \setminus B = [-2; 4) \setminus (1; 6] = [-2; 1]$.
    \end{itemize}
    \item Biểu diễn tập hợp $A \setminus B = [-2; 1]$ trên trục số:
    \begin{center}
    \begin{tikzpicture}[>=Stealth, line width=1.0pt]
        \draw[->] (-4,0) -- (4,0);
        \draw[pattern=north east lines, pattern color=gray!70] (-4,-0.15) rectangle (-2,0.15);
        \draw[pattern=north east lines, pattern color=gray!70] (1,-0.15) rectangle (3.8,0.15);
        \node[font=\bfseries] at (-2,0) {$[$};
        \node[below=3pt, font=\normalsize\bfseries] at (-2,0) {$-2$};
        \node[font=\bfseries] at (1,0) {$]$};
        \node[below=3pt, font=\normalsize\bfseries] at (1,0) {$1$};
    \end{tikzpicture}
    \end{center}
\end{enumerate}

\item \textbf{Câu 2 (0,5 điểm):}
\begin{enumerate}
    \item Áp dụng định lý côsin trong tam giác $ABC$:
    \[
    BC^2 = AB^2 + AC^2 - 2AB\cdot AC\cdot\cos A = 5^2 + 8^2 - 2\cdot 5\cdot 8\cdot\cos 60^\circ = 25 + 64 - 40 = 49.
    \]
    Suy ra $a = BC = \sqrt{49} = 7$.
    \item Diện tích tam giác $ABC$ là:
    \[
    S = \dfrac{1}{2}AB\cdot AC\cdot\sin A = \dfrac{1}{2}\cdot 5\cdot 8\cdot\sin 60^\circ = 20\cdot\dfrac{\sqrt{3}}{2} = 10\sqrt{3}.
    \]
    Mặt khác, $S = \dfrac{1}{2}a\cdot h_a \implies h_a = \dfrac{2S}{a} = \dfrac{2\cdot 10\sqrt{3}}{7} = \dfrac{20\sqrt{3}}{7}$.
\end{enumerate}

\item \textbf{Câu 3 (1,0 điểm):}
Xét tam giác $TAB$:
\begin{itemize}
    \item Góc kề bù với $\widehat{TAH}$ là $\widehat{TAB} = 180^\circ - 40^\circ = 140^\circ$.
    \item Trong tam giác $TAB$, ta có:
    \[
    \widehat{ATB} = 180^\circ - (\widehat{TAB} + \widehat{ABT}) = 180^\circ - (140^\circ + 25^\circ) = 15^\circ.
    \]
    \item Áp dụng định lý sin trong tam giác $TAB$:
    \[
    \dfrac{TA}{\sin\widehat{ABT}} = \dfrac{AB}{\sin\widehat{ATB}} \implies TA = \dfrac{AB\cdot\sin 25^\circ}{\sin 15^\circ} = \dfrac{30\cdot\sin 25^\circ}{\sin 15^\circ} \approx 48{,}977\text{ m}.
    \]
\end{itemize}
Xét tam giác vuông $THA$ tại $H$:
\[
h = TH = TA\cdot\sin\widehat{TAH} = TA\cdot\sin 40^\circ = \dfrac{30\cdot\sin 25^\circ\cdot\sin 40^\circ}{\sin 15^\circ} \approx 31{,}48\text{ m}.
\]
Vậy chiều cao của tháp truyền hình xấp xỉ bằng $31{,}49\text{ m}$.

\item \textbf{Câu 4 (1,0 điểm):}
\begin{enumerate}
    \item Số lượng bò sữa là $x$, số lượng dê là $y$ ($x, y \in \mathbb{N}$).
    \begin{itemize}
        \item Ràng buộc diện tích chuồng: $3x + y \leq 30$.
        \item Ràng buộc vốn tiền thức ăn: $x + 2y \leq 20$.
        \item Điều kiện không âm: $x \geq 0, y \geq 0$.
    \end{itemize}
    Hệ bất phương trình mô tả bài toán là:
    \[
    \begin{cases}
    x \geq 0,\ y \geq 0 \\
    3x + y \leq 30 \\
    x + 2y \leq 20
    \end{cases}
    \]
    \item Miền nghiệm của hệ bất phương trình là miền tứ giác $OABC$ (kể cả bờ) với tọa độ các đỉnh:
    \begin{itemize}
        \item Gốc tọa độ $O(0; 0)$.
        \item Điểm $A(10; 0)$ là giao điểm của đường thẳng $3x + y = 30$ với trục hoành $Ox$.
        \item Điểm $B(8; 6)$ là nghiệm của hệ phương trình $\begin{cases} 3x + y = 30 \\ x + 2y = 20 \end{cases}$.
        \item Điểm $C(0; 10)$ là giao điểm của đường thẳng $x + 2y = 20$ với trục tung $Oy$.
    \end{itemize}
    \item Tổng lợi nhuận thu được sau một đợt nuôi là hàm số:
    \[
    L(x, y) = 4x + 3y\quad (\text{triệu đồng}).
    \]
    Tính giá trị của hàm $L(x, y)$ tại các đỉnh của miền tứ giác:
    \begin{align*}
    L(0; 0) &= 4(0) + 3(0) = 0\text{ (triệu đồng)},\\
    L(10; 0) &= 4(10) + 3(0) = 40\text{ (triệu đồng)},\\
    L(8; 6) &= 4(8) + 3(6) = 32 + 18 = 50\text{ (triệu đồng)},\\
    L(0; 10) &= 4(0) + 3(10) = 30\text{ (triệu đồng)}.
    \end{align*}
    So sánh các giá trị trên, ta thấy $L(x, y)$ đạt giá trị lớn nhất bằng $50\text{ triệu đồng}$ tại đỉnh $B(8; 6)$.
    \par
    \textbf{Kết luận:} Để thu được lợi nhuận lớn nhất, bác Nam nên nuôi $8$ con bò sữa và $6$ con dê. Lợi nhuận lớn nhất thu được là $50\text{ triệu đồng}$.
\end{enumerate}
\end{itemize}

\newpage
\section*{HÌNH VẼ MIỀN NGHIỆM HỆ BẤT PHƯƠNG TRÌNH (BẢN VẼ LỚN IN ẤN SẮC NÉT)}

\begin{center}
\textbf{\large 1. Phần III -- Câu 2: Miền nghiệm tứ giác $OABC$ của hệ bất phương trình}\\[6pt]
\begin{tikzpicture}[scale=1.1, font=\normalsize, >=Stealth, line join=round, line cap=round]
    % Trục toạ độ
    \draw[->, line width=1.2pt] (-1.0,0) -- (7.5,0) node[below=2pt, font=\normalsize\bfseries] {$x$};
    \draw[->, line width=1.2pt] (0,-1.0) -- (0,8.5) node[left=2pt, font=\normalsize\bfseries] {$y$};
    \node[below left=2pt, font=\normalsize\bfseries] at (0,0) {$O$};
    
    % Miền tứ giác O(0,0) - A(4,0) - B(2,4) - C(0,6)
    \fill[ink!15] (0,0) -- (4,0) -- (2,4) -- (0,6) -- cycle;
    
    % Đường 2x + y = 8 (qua (4;0), (2;4), (0;8))
    \draw[line width=1.3pt, blue!80!black] (-0.3,8.6) -- (4.6,-1.2);
    \node[above right, font=\normalsize\bfseries, blue!80!black] at (2.8,3.2) {$d_1\colon 2x+y=8$};
    
    % Đường x + y = 6 (qua (6;0), (2;4), (0;6))
    \draw[line width=1.3pt, teal!80!black] (-0.8,6.8) -- (6.8,-0.8);
    \node[above right, font=\normalsize\bfseries, teal!80!black] at (0.2,6.5) {$d_2\colon x+y=6$};
    
    % Viền cạnh miền tứ giác D
    \draw[line width=2.0pt, ink] (0,0) -- (4,0) -- (2,4) -- (0,6) -- cycle;
    
    % Các đỉnh
    \fill (0,0) circle (3.8pt);
    \fill (4,0) circle (3.8pt);
    \node[below=3pt, font=\normalsize\bfseries, black] at (4,0) {$A(4;0)$};
    \fill (2,4) circle (4.2pt) node[above right=3pt, font=\normalsize\bfseries, red!85!black] {$\mathbf{B(2;4)\ [Max]}$};
    \fill (0,6) circle (3.8pt) node[left=4pt, font=\normalsize\bfseries, black] {$C(0;6)$};
    
    % Đường dóng đỉnh B(2;4)
    \draw[dashed, line width=1.0pt, red!80!black] (2,0) -- (2,4) -- (0,4);
    \node[below=3pt, font=\normalsize\bfseries, black] at (2,0) {$2$};
    \node[left=4pt, font=\normalsize\bfseries, black] at (0,4) {$4$};
    
    % Nhãn miền D
    \node[ink, font=\Large\bfseries] at (1.2, 2.2) {Miền $D$};
\end{tikzpicture}
\end{center}

\vspace{10pt}
\hrule
\vspace{10pt}

\begin{center}
\textbf{\large 2. Phần IV -- Câu 4: Miền nghiệm tứ giác $OABC$ của bài toán tối ưu chăn nuôi}\\[6pt]
\begin{tikzpicture}[scale=0.9, font=\normalsize, >=Stealth, line join=round, line cap=round]
    % Trục toạ độ
    \draw[->, line width=1.2pt] (-1.0,0) -- (12.5,0) node[below=2pt, font=\normalsize\bfseries] {$x$ (bò sữa)};
    \draw[->, line width=1.2pt] (0,-1.0) -- (0,12.5) node[left=2pt, font=\normalsize\bfseries] {$y$ (dê)};
    \node[below left=2pt, font=\normalsize\bfseries] at (0,0) {$O$};
    
    % Miền nghiệm tứ giác O(0,0) - A(10,0) - B(8,6) - C(0,10)
    \fill[ink!15] (0,0) -- (10,0) -- (8,6) -- (0,10) -- cycle;
    
    % Đường 3x + y = 30 (qua (10;0), (8;6), (0;30))
    \draw[line width=1.3pt, blue!80!black] (6.5,10.5) -- (10.6,-1.8);
    \node[above right, font=\normalsize\bfseries, blue!80!black] at (9.0,3.5) {$3x+y=30$};
    
    % Đường x + 2y = 20 (qua (0;10), (8;6), (20;0))
    \draw[line width=1.3pt, teal!80!black] (-0.8,10.4) -- (12.5,3.75);
    \node[above right, font=\normalsize\bfseries, teal!80!black] at (0.2,10.3) {$x+2y=20$};
    
    % Viền miền tứ giác
    \draw[line width=2.0pt, ink] (0,0) -- (10,0) -- (8,6) -- (0,10) -- cycle;
    
    % Các đỉnh
    \fill (0,0) circle (3.8pt);
    \fill (10,0) circle (4.0pt) node[below=3pt, font=\normalsize\bfseries, black] {$A(10;0)$};
    \fill (8,6) circle (4.5pt) node[above right=3pt, font=\normalsize\bfseries, red!85!black] {$\mathbf{B(8;6)\ [Max = 50\text{ tr}]}$};
    \fill (0,10) circle (4.0pt) node[left=4pt, font=\normalsize\bfseries, black] {$C(0;10)$};
    
    % Đường dóng đỉnh B(8;6)
    \draw[dashed, line width=1.0pt, red!80!black] (8,0) -- (8,6) -- (0,6);
    \node[below=3pt, font=\normalsize\bfseries, black] at (8,0) {$8$};
    \node[left=4pt, font=\normalsize\bfseries, black] at (0,6) {$6$};
    
    % Nhãn miền nghiệm
    \node[ink, font=\Large\bfseries] at (4.0, 3.5) {Miền $OABC$};
\end{tikzpicture}
\end{center}
\end{document}
